% Part: first-order-logic
% Chapter: completeness
% Section: introduction

\documentclass[../../../include/open-logic-section]{subfiles}

\begin{document}

\iftag{FOL}
      {\olfileid{fol}{com}{int}}
      {\olfileid{pl}{com}{int}}

\olsection{Inleiding}

De volledigheidsstelling is een van de fundamenteelste resultaten over
logica. Zij heeft twee formuleringen, waarvan we de equivalentie zullen
bewijzen. In haar eerste formulering zegt zij iets fundamenteels over
de verhouding tussen semantisch gevolg en ons !!{derivation}ssysteem:
als !!a{sentence}~$!A$ uit enkele !!{sentence}s $\Gamma$ volgt, dan
bestaat er ook !!a{derivation} die $\Gamma \Proves !A$ vaststelt. Het
!!{derivation}ssysteem is dus zo sterk als het maar kan zijn zonder
dingen te bewijzen die niet werkelijk volgen.

In haar tweede formulering kan zij als een bestaansresultaat voor
modellen worden gesteld: elke consistente verzameling !!{sentence}s is
vervulbaar. Consistentie is een bewijstheoretisch begrip: het zegt dat
ons !!{derivation}ssysteem bepaalde !!{derivation}s niet kan
voortbrengen. Maar waarom zou alleen het ontbreken van een bepaald
soort !!{derivation}s uit~$\Gamma$ garanderen dat er
\iftag{FOL}{!!a{structure}~$\Struct{M}$}{!!{valuation}{~$\pAssign{v}$}
bestaat waarvoor $\iftag{FOL}{\Sat{M}}{\pSat{v}}{\Gamma}$}? Voordat de
volledigheidsstelling voor het eerst werd bewezen---zelfs voordat de
!!{derivation}ssystemen bestonden die we nu gebruiken---was de grote
Duitse wiskundige David Hilbert van mening dat de consistentie van
wiskundige theorieën het bestaan garandeert van de objecten waarover
ze gaan. In een brief aan Gottlob Frege formuleerde hij dit als volgt:
\begin{quote}
  Als de willekeurig gegeven axioma's elkaar met al hun gevolgen niet
  tegenspreken, dan zijn ze waar en bestaan de door de axioma's
  gedefinieerde dingen. Dit is voor mij het criterium voor waarheid en
  bestaan.
\end{quote}
Frege was het daar hartgrondig mee oneens. De tweede formulering van de
volledigheidsstelling laat zien dat Hilbert ten minste in deze zin
gelijk had: als de axioma's consistent zijn, bestaat er \emph{een}
\iftag{FOL}{!!{structure}}{!!{valuation}} die ze allemaal waar maakt.

Dit zijn niet de enige redenen waarom de volledigheidsstelling---of
eigenlijk haar bewijs---belangrijk is. Zij heeft verscheidene
belangrijke gevolgen, waarvan we sommige afzonderlijk zullen
bespreken. Elke !!{derivation} die $\Gamma \Proves !A$ aantoont, is
bijvoorbeeld eindig en kan dus slechts eindig veel !!{sentence}s
uit~$\Gamma$ gebruiken. Uit de volledigheidsstelling volgt daarom dat,
als $!A$ een gevolg is van~$\Gamma$, zij al een gevolg is van een
eindige deelverzameling van~$\Gamma$. Dit heet \emph{compactheid}.
Equivalent hiermee is: als elke eindige deelverzameling van $\Gamma$
consistent is, dan moet $\Gamma$ zelf consistent zijn.

Hoewel de compactheidsstelling via de omweg langs !!{derivation}s uit
de volledigheidsstelling volgt, kan men ook \emph{het bewijs van} de
volledigheidsstelling gebruiken om haar rechtstreeks aan te tonen. Het
bewijs neemt namelijk een verzameling !!{sentence}s met een bepaalde
eigenschap---consistentie---en construeert daaruit !!a{structure} met
bepaalde eigenschappen (in dit geval een structuur die de verzameling
vervult). Nagenoeg dezelfde constructie bewijst compactheid
rechtstreeks wanneer men niet met consistente, maar met ``eindig
vervulbare'' verzamelingen !!{sentence}s begint.
\iftag{FOL}{De constructie heeft nog andere gevolgen. Elke vervulbare
  verzameling !!{sentence}s heeft bijvoorbeeld een eindig of
  !!{denumerable} model. (Dit resultaat heet de stelling van
  L\"owenheim--Skolem.) In de logica, en soms zelfs in de filosofie,
  construeert men in het algemeen vaak !!{structure}s uit verzamelingen
  !!{sentence}s.}{}

\end{document}
